\usepackage{listings}

\renewcommand{\lstlistlistingname}{Lista de Programas}
\renewcommand{\lstlistingname}{Programa}

\DeclareFixedFont{\ttb}{T1}{txtt}{bx}{n}{10} % for bold
\DeclareFixedFont{\ttm}{T1}{txtt}{m}{n}{10}  % for normal

\newcommand{\gxnumbersep}{0.5em}

% um set default para listagens
\lstset{extendedchars=true,%
    basicstyle=\ttfamily,%
    inputencoding=utf8,%
    literate=%
    {ã}{{\~{a}}}1%
    {á}{{\'{a}}}1%
    {â}{{\^{a}}}1%
    {à}{{\`{a}}}1%
    {é}{{\'{e}}}1%
    {è}{{\`{e}}}1%
    {ê}{{\^{e}}}1%
    {î}{{\^{i}}}1%
    {í}{{\'{i}}}1%
    {ô}{{\^{o}}}1%
    {ó}{{\'{o}}}1%
    {õ}{{\~{o}}}1%
    {û}{{\^{u}}}1%
    {ú}{{\'{u}}}1%
    {ç}{{\c{c}}}1%
    {Ç}{{\c{C}}}1%
    {À}{{\`{A}}}1%
    {Â}{{\^{A}}}1%
    {Á}{{\'{A}}}1%
    {É}{{\'{E}}}1%
    {Ê}{{\^{E}}}1%
    {Í}{{\'{I}}}1%
    {Ó}{{\'{O}}}1%
    {Ô}{{\^{O}}}1%
    {Õ}{{\~{O}}}1%
    {Ú}{{\'{U}}}1%
}
%   https://tex.stackexchange.com/questions/83882/how-to-highlight-python-syntax-in-latex-listings-lstinputlistings-command
\definecolor{deepblue}{rgb}{0,0,0.5}
\definecolor{deepred}{rgb}{0.6,0,0}
\definecolor{deepgreen}{rgb}{0,0.5,0}

\newcommand\pythonstyle{\lstset{%
        language=Python,%
        basicstyle=\ttm,%
        %basicstyle=\scriptsize,
        numbers=left,%
        numberstyle=\small\ttfamily, %
        numbersep=\gxnumbersep,%
        commentstyle=\color{purple},
        morekeywords={self,as},              % Add keywords here
        keywordstyle=\ttb\color{deepblue},%
        emph={MyClass,__init__},          % Custom highlighting
        emphstyle=\ttb\color{deepred},    % Custom highlighting style
        stringstyle=\color{deepgreen},%
        frame=tb,                         % Any extra options here
        showstringspaces=false,%
        extendedchars=true,%
        firstnumber = auto,%
        inputencoding=utf8,%
        breaklines=true,%
        escapechar=`,%
        postbreak=\mbox{\textcolor{red}{$\hookrightarrow$}\space}
 }}


%
\makeatletter
    \lstnewenvironment{python}[1][]
    {   \pythonstyle
        \lstset{#1}
        \csname\@lst @SetFirstNumber\endcsname
    }
    {     \csname \@lst @SaveFirstNumber\endcsname
    }
\makeatother
    % Python for inline
    \newcommand\pythoninline[1]{{\pythonstyle\lstinline!#1!}}
%
% Python for external files
\newcommand\pythonexternal[2][]{{
        \pythonstyle
        \lstinputlisting[#1]{#2}}}


\newcommand\xmlstyle{\lstset{%
        language=XML,%
        basicstyle=\ttm,%
        %basicstyle=\scriptsize,
        numbers=left,%
        numberstyle=\small\ttfamily, %
        numbersep=\gxnumbersep,%
        morekeywords={xs:schema,xs:element,xs:sequence,xs:complexType},              % Add keywords here
        keywordstyle=\ttb\color{deepblue},%
        emph={name,type,attributeFormDefault,elementFormDefault,
        xmlns:xs},          % Custom highlighting
        emphstyle=\ttb\color{deepred},    % Custom highlighting style
        stringstyle=\color{deepgreen},%
        frame=tb,                         % Any extra options here
        showstringspaces=false,%
        extendedchars=true,%
        inputencoding=utf8,%
        breaklines=true,%
        postbreak=\mbox{\textcolor{red}{$\hookrightarrow$}\space}
}}


  \lstnewenvironment{xml}[1][]
{   \xmlstyle
    \lstset{#1}
}
{}

\newcommand\htmlstyle{\lstset{%
        language=HTML,%
        basicstyle=\ttm,%
        %basicstyle=\scriptsize,
        numbers=left,%
        numberstyle=\small\ttfamily, %
        numbersep=\gxnumbersep,%
        %morekeywords={},              % Add keywords here
        keywordstyle=\ttb\color{deepblue},%
        %emph={},          % Custom highlighting
        emphstyle=\ttb\color{deepred},    % Custom highlighting style
        stringstyle=\color{deepgreen},%
        frame=tb,                         % Any extra options here
        showstringspaces=false,%
        extendedchars=true,%
        inputencoding=utf8,%
        breaklines=true,%
        postbreak=\mbox{\textcolor{red}{$\hookrightarrow$}\space}
}}


\lstnewenvironment{html}[1][]
{   \htmlstyle
    \lstset{#1}
}
{}



\DefineVerbatimEnvironment{gxoutput}{Verbatim}{%
frame=lines,numbers=left,numbersep=\gxnumbersep}
\renewcommand{\theFancyVerbLine}{%
{\small\ttfamily \arabic{FancyVerbLine}}}

\DefineVerbatimEnvironment{gxoutputs}{Verbatim}{%
frame=lines,numbers=left,numbersep=\gxnumbersep,
fontsize=\small}
\renewcommand{\theFancyVerbLine}{%
{\small\ttfamily \arabic{FancyVerbLine}}}





\newcommand\prologstyle{\lstset{%
        language=Prolog,%
        basicstyle=\ttm,%
        %basicstyle=\scriptsize,
        numbers=left,%
        numberstyle=\small\ttfamily, %
        numbersep=\gxnumbersep,%
        commentstyle=\color{purple},
%        morekeywords={self,as},              % Add keywords here
        keywordstyle=\ttb\color{deepblue},%
        emph={MyClass,__init__},          % Custom highlighting
        emphstyle=\ttb\color{deepred},    % Custom highlighting style
        stringstyle=\color{deepgreen},%
        frame=tb,                         % Any extra options here
        showstringspaces=false,%
        extendedchars=true,%
        firstnumber = auto,%
        inputencoding=utf8,%
        breaklines=true,%
        escapechar=`,%
        postbreak=\mbox{\textcolor{red}{$\hookrightarrow$}\space}
 }}

\makeatletter
    \lstnewenvironment{prolog}[1][]
    {   \prologstyle
        \lstset{#1}
        \csname\@lst @SetFirstNumber\endcsname
    }
    {     \csname \@lst @SaveFirstNumber\endcsname
    }
\makeatother



\newcommand\javastyle{\lstset{%
		language=Java,%
		basicstyle=\ttm,%
		%basicstyle=\scriptsize,
		numbers=left,%
		numberstyle=\small\ttfamily, %
		numbersep=\gxnumbersep,%
		commentstyle=\color{purple},
		morekeywords={self,as},              % Add keywords here
		keywordstyle=\ttb\color{deepblue},%
%		emph={MyClass,__init__},          % Custom highlighting
		emphstyle=\ttb\color{deepred},    % Custom highlighting style
		stringstyle=\color{deepgreen},%
		frame=tb,                         % Any extra options here
		showstringspaces=false,%
		extendedchars=true,%
		firstnumber = auto,%
		inputencoding=utf8,%
		breaklines=true,%
		escapechar=`,%
		postbreak=\mbox{\textcolor{red}{$\hookrightarrow$}\space}
}}


%
\makeatletter
\lstnewenvironment{java}[1][]
{   \javastyle
	\lstset{#1}
	\csname\@lst @SetFirstNumber\endcsname
}
{     \csname \@lst @SaveFirstNumber\endcsname
}
\makeatother
% java for inline
\newcommand\javainline[1]{{\javastyle\lstinline!#1!}}
%
% java for external files
\newcommand\javaexternal[2][]{{
		\javastyle
		\lstinputlisting[#1]{#2}}}

% ALGORITHM 2E Settings

\usepackage[portuguese,linesnumbered,lined,algochapter,ruled ]{algorithm2e}
\SetKwFor{Para}{para}{faça}{fim do para}
\SetKwFor{Enquanto}{enquanto}{faça}{fim do enquanto}

